\documentclass[11pt,letterpaper]{article}

\usepackage[T1]{fontenc}
\usepackage{amsmath,amssymb}
\usepackage{tikz}
\usepackage{pgfplots}
\pgfplotsset{compat=1.18}
\usepackage{algorithm}
\usepackage{algpseudocode}
\usepackage{booktabs}
\usepackage{enumitem}
\usepackage[hidelinks]{hyperref}

\newcommand{\N}{\mathbb{N}}
\newcommand{\Z}{\mathbb{Z}}
\newcommand{\Q}{\mathbb{Q}}
\newcommand{\R}{\mathbb{R}}
\newcommand{\C}{\mathbb{C}}
\newcommand{\Map}[3]{#1 \colon #2 \to #3}
\newcommand{\Mod}[1]{#1\text{-}\mathrm{Mod}}
\newcommand{\GroupHom}{\mathrm{Hom}_{\mathbf{Grp}}}
\newcommand{\RingHom}{\mathrm{Hom}_{\mathbf{Ring}}}

\newtheorem{theorem}{Theorem}
\newtheorem{lemma}[theorem]{Lemma}
\newtheorem{corollary}[theorem]{Corollary}
\newtheorem{definition}[theorem]{Definition}
\newtheorem{remark}[theorem]{Remark}
\newenvironment{proof}{\noindent\textit{Proof.}}{\hfill\(\square\)}

\title{No uniform bound on one affine class}
\author{ascendnoosphere}
\date{\today}

\begin{document}
\maketitle

\begin{abstract}
If a time \(j\) is the first time the divisor of a parameter record divides the
affine state of a seed, and if the multiplier power at that time is at least
the absolute value of the divisor, then no single pair of natural numbers
\(A<|d|\) and \(B\) bounds \(|d|\) times the absolute value of the quotient
by \(A|n|+B\) for every integer in the residue class of the seed. The argument
moves the seed by natural multiples of \(|d|\), keeps the first hit, and
compares one seed whose absolute value exceeds an explicit threshold. A single
seed in the class may still satisfy the inequality.
\end{abstract}

\section{Affine state and first hit}

Let \(q,c\in\Z\). Define \(T_0(q)=0\) and \(T_{j+1}(q)=T_j(q)+q^j\) for every
natural number \(j\). Thus \(T_j(q)=\sum_{i=0}^{j-1}q^i\), with the empty sum
equal to zero.

\begin{definition}
For \(q,c,n\in\Z\) and \(j\in\N\), the affine state is
\[
A_j(q,c,n)=q^j n+c\,T_j(q).
\]
\end{definition}

A parameter record \(p\) consists of integers \(q\), \(c\), and \(d\) with
\(|d|\ge 2\). In particular \(d\neq 0\). Write \(A_j(n)\) for
\(A_j(q,c,n)\) when the record is fixed.

\begin{definition}
The time \(j\in\N\) is a first hit of the seed \(n\in\Z\), written
\(\operatorname{Hit}(p,n,j)\), when \(d\) divides \(A_j(n)\) and \(d\) divides
\(A_i(n)\) for no natural number \(i<j\).
\end{definition}

The residue class of \(n_0\) is taken modulo \(|d|\). Two integers \(n\) and
\(n_0\) lie in the same class when their images in \(\Z/|d|\Z\) agree, or
equivalently when \(|d|\) divides \(n-n_0\).

\section{Transport by a multiple of the divisor}

\begin{lemma}
\label{lem:shift}
For all \(q,c,n,t,d\in\Z\) and every \(j\in\N\),
\[
A_j(q,c,n+td)=A_j(q,c,n)+q^j td.
\]
\end{lemma}

\begin{proof}
Expand both sides. The offset sum depends on \(q\) and \(j\) only, so
\[
A_j(q,c,n+td)=q^j(n+td)+c\,T_j(q)
=q^j n+c\,T_j(q)+q^j td.
\]
The first two summands are \(A_j(q,c,n)\).
\end{proof}

\begin{lemma}
\label{lem:difference}
For all \(q,c,n,t,d\in\Z\) and every \(j\in\N\), the integer \(d\) divides
\(A_j(q,c,n+td)-A_j(q,c,n)\).
\end{lemma}

\begin{proof}
Lemma~\ref{lem:shift} gives the difference \(q^j td\), which is
\(d\cdot(q^j t)\).
\end{proof}

\begin{lemma}
\label{lem:div-shift}
Let \(p\) be a parameter record with divisor \(d\). For every \(n,t\in\Z\) and
every \(i\in\N\),
\[
d\mid A_i(n)\quad\text{if and only if}\quad d\mid A_i(n+td).
\]
\end{lemma}

\begin{proof}
Let \(k\) be the integer given by Lemma~\ref{lem:difference}, so
\(A_i(n+td)-A_i(n)=dk\). If \(A_i(n)=dm\), then
\(A_i(n+td)=d(m+k)\). If \(A_i(n+td)=dm\), then
\(A_i(n)=d(m-k)\).
\end{proof}

\begin{lemma}
\label{lem:hit-shift}
Let \(p\) be a parameter record with divisor \(d\). For every \(n,t\in\Z\) and
every \(j\in\N\),
\[
\operatorname{Hit}(p,n,j)\quad\text{if and only if}\quad
\operatorname{Hit}(p,n+td,j).
\]
\end{lemma}

\begin{proof}
Lemma~\ref{lem:div-shift} is an equivalence at time \(j\) and at every
\(i<j\). The divisibility clause and the failure clause therefore transport
together, in both directions.
\end{proof}

\begin{lemma}
\label{lem:sign}
For every \(d\in\Z\),
\[
|d|=
\begin{cases}
d & \text{if }0\le d,\\
-d & \text{if }d<0.
\end{cases}
\]
Equivalently, \(|d|=\varepsilon d\) in \(\Z\), where \(\varepsilon=1\) if
\(0\le d\) and \(\varepsilon=-1\) if \(d<0\).
\end{lemma}

\begin{proof}
If \(0\le d\), then \(|d|=d=1\cdot d\). If \(d<0\), then \(0\le -d\) and
\(|d|=|-d|=-d=(-1)\cdot d\).
\end{proof}

\begin{lemma}
\label{lem:residue}
For every \(n_0\in\Z\) and all \(m,t\in\N\),
\[
n_0+tm\equiv n_0\pmod{m}
\]
in \(\Z/m\Z\). If \(m=0\), both sides are interpreted in the zero ring, and
the multiple of \(m\) is still zero.
\end{lemma}

\begin{proof}
The image of \(m\) in \(\Z/m\Z\) is zero, so the image of \(tm\) is zero.
Adding it does not change the image of \(n_0\).
\end{proof}

\begin{lemma}
\label{lem:natabs-shift}
Let \(p\) be a parameter record with divisor \(d\), and let \(n_0\in\Z\) and
\(j,t\in\N\). If \(\operatorname{Hit}(p,n_0,j)\), then
\[
\operatorname{Hit}(p,\,n_0+t|d|,\,j).
\]
\end{lemma}

\begin{proof}
Let \(\varepsilon\) be the sign from Lemma~\ref{lem:sign}, so \(|d|=\varepsilon d\)
as integers. Then
\[
t|d|=t(\varepsilon d)=(t\varepsilon)d=(\varepsilon t)d.
\]
Lemma~\ref{lem:hit-shift}, applied to the integer shift \(\varepsilon t\),
gives the first hit of \(n_0+(\varepsilon t)d\).
\end{proof}

The sequence \(n_0+t|d|\) is the explicit increasing sequence of absolute
values used below. It is not necessary to invoke infinitude of the class
separately: the next lemma names a particular \(t\).

\begin{lemma}
\label{lem:exists-t}
Let \(n_0\in\Z\) and let \(m,B\in\N\) with \(m\ge 1\). There exists \(t\in\N\)
such that
\[
B\le |n_0+tm|.
\]
One such value is \(t=|n_0|+B\).
\end{lemma}

\begin{proof}
Set \(t=|n_0|+B\). Since \(m\ge 1\) and \(t\ge 0\),
\[
t=t\cdot 1\le t m.
\]
Hence \(|n_0|+B\le tm\). Also \(tm\ge 0\), so \(|tm|=tm\). The splitting
\(tm=(n_0+tm)+(-n_0)\) and the triangle inequality give
\[
tm=|tm|\le |n_0+tm|+|n_0|.
\]
Therefore
\[
|n_0|+B\le |n_0+tm|+|n_0|.
\]
Cancel the nonnegative integer \(|n_0|\) on the right to obtain
\(B\le |n_0+tm|\).
\end{proof}

\section{Quotient and reverse triangle}

\begin{lemma}
\label{lem:product}
Let \(d,a\in\Z\) and suppose \(d\) divides \(a\). Then
\[
|d|\cdot |a/d|=|a|.
\]
\end{lemma}

\begin{proof}
Exact division gives \((a/d)\cdot d=a\). Absolute value is multiplicative, so
\[
|a/d|\cdot |d|=|a|.
\]
Commutativity of multiplication of natural numbers swaps the factors.
\end{proof}

\begin{lemma}
\label{lem:reverse}
Let \(q,c,n\in\Z\) and \(j\in\N\). If
\[
|c|\,|T_j(q)|\le |q|^j\,|n|,
\]
then
\[
|q|^j\,|n|-|c|\,|T_j(q)|\le |A_j(q,c,n)|.
\]
The subtraction on the left is the natural-number subtraction. The hypothesis
makes it exact.
\end{lemma}

\begin{proof}
The linear term splits as
\[
q^j n=A_j(q,c,n)+(-(c\,T_j(q))).
\]
The triangle inequality and \(|-x|=|x|\) give
\[
|q^j n|\le |A_j(q,c,n)|+|c\,T_j(q)|.
\]
Subtract \(|c\,T_j(q)|\) from both sides in \(\Z\). Multiplicativity of
absolute value converts the three terms into the natural numbers
\(|q|^j|n|\), \(|c|\,|T_j(q)|\), and \(|A_j(q,c,n)|\). The hypothesis lets the
integer subtraction of the first two pass to the corresponding natural-number
subtraction.
\end{proof}

\section{One large seed}

\begin{lemma}
\label{lem:large}
Let \(p\) be a parameter record with multiplier \(q\), offset \(c\), and
divisor \(d\). Let \(n\in\Z\), let \(j,A,B\in\N\), and assume
\(\operatorname{Hit}(p,n,j)\), \(|d|\le |q|^j\), and \(A<|d|\). Set
\(M=|q|^j\) and \(K=|c|\,|T_j(q)|\). If
\[
B+K+1\le |n|
\]
and
\[
|d|\cdot |A_j(n)/d|\le A|n|+B,
\]
then the two displayed inequalities cannot hold simultaneously.
\end{lemma}

\begin{proof}
The first-hit hypothesis supplies \(d\mid A_j(n)\). Lemma~\ref{lem:product}
converts the proposed bound into
\[
|A_j(n)|\le A|n|+B.
\]
Since \(A<|d|\) and \(|d|\le M\), one has \(A+1\le M\), hence \(A\le M\) and
\[
1\le M-A.
\]
The size hypothesis gives \(K\le B+K+1\le |n|\). Also \(M\ge |d|\ge 2\), so
\(1\le M\) and \(|n|\le M|n|\). Therefore \(K\le M|n|\).
Lemma~\ref{lem:reverse} yields
\[
M|n|-K\le |A_j(n)|\le A|n|+B.
\]
Adding \(K\) to both sides of the outer comparison, using that the natural
subtraction is exact under \(K\le M|n|\), gives
\[
M|n|\le A|n|+(B+K),
\]
and then
\[
M|n|-A|n|\le B+K.
\]
The left side is \((M-A)|n|\) because \(A\le M\). Since \(1\le M-A\),
\[
|n|\le (M-A)|n|\le B+K.
\]
This contradicts \(B+K+1\le |n|\).
\end{proof}

\begin{theorem}
\label{thm:no-uniform}
Let \(p\) be a parameter record with multiplier \(q\), offset \(c\), and
divisor \(d\). Let \(n_0\in\Z\) and \(j\in\N\). Assume
\(\operatorname{Hit}(p,n_0,j)\) and \(|d|\le |q|^j\). There do not exist
\(A,B\in\N\) with \(A<|d|\) such that for every \(n\in\Z\) in the residue
class of \(n_0\) modulo \(|d|\),
\[
|d|\cdot |A_j(n)/d|\le A|n|+B.
\]
\end{theorem}

\begin{proof}
Suppose \(A\) and \(B\) were such natural numbers. Set
\(K=|c|\,|T_j(q)|\) and \(m=|d|\). Then \(m\ge 2\ge 1\).
Lemma~\ref{lem:exists-t} supplies a natural number \(t\) with
\[
B+K+1\le |n_0+tm|.
\]
Set \(n=n_0+t|d|\). Lemma~\ref{lem:residue} puts \(n\) in the class of \(n_0\),
so the supposed bound applies to \(n\). Lemma~\ref{lem:natabs-shift} gives
\(\operatorname{Hit}(p,n,j)\). The power hypothesis is unchanged, and
\(A<|d|\). Lemma~\ref{lem:large} then says that the bound and the size
inequality cannot both hold. This contradiction shows that no such pair
\(A,B\) exists.
\end{proof}

The quantifier is on the pair. Theorem~\ref{thm:no-uniform} does not say that
every seed fails every inequality of this shape. A seed with small absolute
value can satisfy \(A|n|+B\) for a pair that already fails on a larger seed of
the same class.

\begin{remark}
\label{rem:example}
Take \(q=2\), \(c=4\), \(d=4\), \(n_0=1\), and \(j=2\). Then \(|d|=4=|q|^j\),
\(T_2(2)=1+2=3\), and
\[
A_2(2,4,n)=4n+12.
\]
For \(n_0=1\), the state is \(16\), while the earlier states are \(1\) and
\(6\). Thus \(j=2\) is a first hit. The seed \(n=1\) satisfies
\[
|d|\cdot |A_2(1)/d|=16\le 3\cdot 1+1000.
\]
The seed \(n=1001=1+250\cdot 4\) is in the same class, and the states at times
\(0\) and \(1\) are \(1001\) and \(2006\), neither divisible by \(4\), while
\[
A_2(2,4,1001)=4016,\qquad 4016/4=1004.
\]
Hence \(4\cdot 1004=4016>3\cdot 1001+1000=4003\). The same first hit and the
same absolute quotient occur for the divisor \(-4\), because a sign change of
the divisor does not change divisibility and replaces the quotient by its
negative. The pair \(A=3\), \(B=1000\) therefore fails on this class, although
it holds at the original seed.
\end{remark}

\begin{theorem}
\label{thm:trichotomy}
Fix a seed \(n_0\). Exactly one of the following descriptions applies to the
integers \(n\equiv n_0\pmod{|d|}\). No such \(n\) has a first hit. Or there
is one time \(j\) such that every such \(n\) has its first hit at \(j\) and
\(|d|\le |q|^j\). Or there is one time \(j\) such that every such \(n\) has
its first hit at \(j\) and \(|q|^j<|d|\).
\end{theorem}

\begin{proof}
If \(|d|\) divides \(n'-n\), write \(|d|=\varepsilon d\) with
\(\varepsilon=\pm 1\). Then \(n' = n+(\varepsilon k)d\) for the quotient
\(k\), so the first-hit time is the same. Thus either \(n_0\) has no first
hit, in which case no congruent seed has one, or \(n_0\) has a first hit at
some \(j\), in which case every congruent seed has that same first hit. The
power \(|q|^j\) is then either at least \(|d|\) or strictly smaller. These
branches exclude one another. The modulus divides \(0\), so \(n_0\) lies in
its own class: a first hit of \(n_0\) refutes the empty branch. A power cannot
be both strictly smaller than \(|d|\) and at least \(|d|\).
\end{proof}

\end{document}
